\documentclass[11pt, a4paper]{letter}
\usepackage[utf8]{inputenc}
\usepackage[margin=1in]{geometry}
\usepackage{hyperref}

\begin{document}
\begin{letter}{Dr. Benoit Caillaud \\ HYCOMES Team \\ Inria Rennes - Bretagne Atlantique \\ Rennes, France}
\opening{Dear Dr. Benoit Caillaud,}

I am writing to express my strong interest in the PhD position focusing on Structural Methods for Mixed Model/Data Digital Twin Engineering within the HYCOMES team. As an AI \& Data Science Engineer with two rigorous engineering theses dedicated entirely to Digital Twin architectures and physics-informed machine learning, this research aligns perfectly with my background and career ambitions.

In my State Engineer thesis ("Gas Turbine Digital Twin"), I faced the exact challenge of model/data mismatch. Integrating high-frequency SCADA telemetry with theoretical thermodynamic performance models required resolving immense inconsistencies caused by sensor drift and unmodeled mechanical wear. I engineered predictive maintenance architectures using Spatial-Temporal Transformers that learned these discrepancies. However, I quickly realized that purely data-driven optimization does not scale robustly to complex dynamical systems.

This realization drove my Master's thesis on the "Cardiovascular Medical Digital Twin." To ensure rigorous model initialization and fine-tuning, I developed PI-DeepONet, an architecture that explicitly maps the differential-algebraic equations (DAEs) of the 0D Windkessel hemodynamic model into the neural architecture via PyTorch Autograd. I also engineered a "Missingness-Aware Gate" to systematically handle observational data gaps. Transitioning from these data-heavy approaches to structural analysis for DAE systems—specifically computing Minimal Structurally Overdetermined (MSO) subsystems as parity spaces—is the exact methodological leap I wish to make for my doctoral research.

I possess strong programming skills in Python and PyTorch, a rigorous mathematical foundation in dynamical systems, and a deep familiarity with the limitations of current data assimilation techniques in digital twin engineering. 

Please find attached my CV, research statement, and academic transcripts. I would welcome the opportunity to discuss my qualifications with you.

\closing{Sincerely, \\ Ismail Aissaoui \\ \href{https://ismail-aissaoui.vercel.app}{https://ismail-aissaoui.vercel.app}}
\end{letter}
\end{document}